\documentclass{l4dc2026}
\usepackage[T1]{fontenc}
\usepackage{times,booktabs}
\jmlrvolume{}
\jmlrproceedings{ESE 6180}{ESE 6180 Midterm Report}
\jmlrworkshop{Midterm Report}
\title[History and Action Chunking]{History and Action Chunking in Stable Linear Imitation Learning}
\author{\Name{Nihar Ballamudi}\and\Name{Kyle Horton}\\
\addr University of Pennsylvania}
\newcommand{\norm}[1]{\left\lVert#1\right\rVert}
\newcommand{\spr}{\operatorname{spr}}
\begin{document}
\maketitle
\begin{abstract}
We study when history can replace action chunking in imitation of a stabilizing linear expert. We characterize exact action recovery from observation--action histories, establish noiseless equivalence of execution horizons under sufficient context, and derive an incremental stability bound for chunked policies. History resolves missing information; commitment can instead suppress unstable feedback from an imperfect predictor. Supporting simulations illustrate both mechanisms and the separate role of training coverage. The remaining objective is to connect context length to learned-policy error and the shortest reliable execution horizon.
\end{abstract}
\begin{keywords}
Imitation learning, observability, action chunking, incremental stability.
\end{keywords}

\section{Question and mathematical setting}
Restricted demonstrations can leave corrective behavior unidentified even when the plant and expert are stable \citep{simchowitz2025}. Action chunking can exploit stable open-loop dynamics to limit error amplification \citep{zhang2025}. We ask when additional history makes frequent replanning reliable, reducing the need for commitment.

Consider $x_{t+1}=Ax_t+Bu_t+w_t$, $y_t=Cx_t$, and expert $u_t^*=Kx_t$; write $F=A+BK$. Schur stability of $A$ concerns decay under common inputs, whereas stability of $F$ concerns expert feedback. The expert is Markovian in $x$; the learner observes
\[
z_t^{(r)}=\operatorname{col}(y_{t-r+1:t},u_{t-r+1:t-1}).
\]
Context $r$ counts observations, prediction horizon $H$ counts predicted actions, and execution length $\ell\le H$ counts committed actions. At $b=j\ell$, the learner predicts a chunk and applies $u_{b+k}=\pi_k(z_b)$ for $0\le k<\ell$. New measurements do not revise that chunk. Reconstruction statements assume $w=0$, known applied inputs, and a valid initial history. The following results are elementary control derivations for this project, not novelty claims about observability.

\section{History and exact imitation}
Let $\mathcal O_r=\operatorname{col}(C,CA,\ldots,CA^{r-1})$ and $h=\min\{r:\operatorname{rank}\mathcal O_r=n\}$ for an observable pair. With $s=t-r+1$, unrolling gives $Y=\mathcal O_rx_s+\mathcal T_rU$ and $x_t=A^{r-1}x_s+\mathcal R_rU$. Here $Y,U$ stack the observed outputs and applied inputs, $\mathcal R_r=[A^{r-2}B\ \cdots\ B]$, and $(\mathcal T_r)_{ij}=CA^{i-1-j}B$ for $j<i$, zero otherwise (zero-based blocks).

\textbf{Lemma 1 (reconstruction).} If $r\ge h$, then $x_t=R_rz_t$, where
\begin{equation}
R_r=[A^{r-1}\mathcal O_r^\dagger\quad\mathcal R_r-A^{r-1}\mathcal O_r^\dagger\mathcal T_r].
\label{eq:reconstruction}
\end{equation}
\emph{Proof.} Recover $x_s=\mathcal O_r^\dagger(Y-\mathcal T_rU)$ and propagate to $t$. $\square$ The input correction is essential; the identity holds under arbitrary learner actions.

\textbf{Proposition 1 (action recovery).} A policy satisfying $\pi_r(z_t)=Kx_t$ for every noiseless history exists if and only if
\begin{equation}\ker\mathcal O_r\subseteq\ker(KA^{r-1}).\label{eq:kernel}\end{equation}
\emph{Proof.} Under identical inputs, indistinguishable histories have initial-state difference $d\in\ker\mathcal O_r$ and current-action difference $KA^{r-1}d$. Necessity follows. Conversely, kernel inclusion makes $\mathcal O_rx_s\mapsto KA^{r-1}x_s$ a well-defined linear map; add the known input contribution. $\square$

Full observability is sufficient but not necessary: only action-relevant information must be recovered. If \eqref{eq:kernel} fails, indistinguishable histories require actions separated by some $\Delta\ne0$. Every deterministic predictor errs by at least $\norm{\Delta}/2$ on one history. The first action of any chunk obeys the same lower bound. Commitment cannot create missing information. In fact, exact prediction of an entire $\ell$-action expert chunk requires
\[
\ker\mathcal O_r\subseteq\bigcap_{k=0}^{\ell-1}\ker(KF^kA^{r-1}),
\]
since its $k$th action is $KF^kx_t$. Longer prediction can demand more information than the current action alone.

\textbf{Proposition 2 (execution equivalence).} For $r\ge h$, the policy $\pi_k^*(z)=KF^kR_rz$ reproduces the noiseless expert trajectory for every $1\le\ell\le H$. Consequently, if $F$ is Schur stable, exact single-step imitation is stable.

\emph{Proof.} Reconstruction gives the boundary state $x_b$. If $x_{b+k}=F^kx_b$, applying $KF^kx_b$ gives $x_{b+k+1}=F^{k+1}x_b$. Induct within each chunk and repeat at the next boundary. $\square$

Thus sufficient context removes a representational need for chunking, but no horizon is strictly preferable in the exact noiseless case. Learning error and disturbances create the tradeoff. Under shared disturbances, single-step learner--expert tracking error satisfies $d_{t+1}=Fd_t+Be_t$, where $e_t=\hat\pi(z_t)-Kx_t$. If $\norm{F^k}\le M_F\beta^k$, $\beta<1$, and deployment error is bounded by $\varepsilon$, then $\limsup_t\norm{d_t}\le M_F\norm{B}\varepsilon/(1-\beta)$. More explicitly, unrolling gives
\[
\norm{d_t}\le M_F\beta^t\norm{d_0}
+M_F\norm{B}\sum_{i=0}^{t-1}\beta^{t-1-i}\norm{e_i}.
\]
Small training loss does not establish the required deployment bound. If only observations are noisy, reconstruction error is $A^{r-1}\mathcal O_r^\dagger V_t$, where $V_t$ stacks measurement errors. Full rank therefore does not ensure good conditioning or robustness.

\paragraph{Representation versus identification.}
Noiseless expert histories have the form $z_t=S_rx_{t-r+1}$, with $S_r$ stacking $CF^j$ for $0\le j<r$ and $KF^j$ for $0\le j<r-1$. Their feature rank is at most $n$. Any $\Delta W$ satisfying $\Delta WS_r=0$ is invisible on demonstrations but may change behavior on disturbed histories. More context can therefore enlarge the policy class without identifying its off-demonstration behavior. Process noise also invalidates exact deterministic reconstruction.

\clearpage
\section{A stability guarantee for action chunks}
Set $w=0$. Assume that all admissible trajectory pairs satisfy
\begin{equation}
\norm{\delta y_{t+k}}\le C_0\rho^k\norm{\delta z_t}_{\mathcal Z}
+C_0\sum_{i=0}^{k-1}\rho^{k-1-i}\norm{\delta u_{t+i}},
\label{eq:ios}
\end{equation}
where $C_0\ge1$, $0<\rho<1$, and $\norm{z}_{\mathcal Z}$ sums block Euclidean norms. For observable Schur-stable plants and $r\ge h$, this follows from Lemma 1 and $\norm{A^k}\le M_A\rho^k$, with $C_0$ enlarged to bound the output and input maps. Observability alone does not imply decay.

\textbf{Proposition 3 (history-based chunk stability).} Use the same policy and schedule on both trajectories, with $u_{j\ell+k}=\pi_k(z_{j\ell})+v_{j\ell+k}$ and $\operatorname{Lip}_{\mathcal Z}(\pi_k)\le D\rho^k$. For $0<\alpha<1$, define
\[
\bar\rho=\rho^{1-\alpha},\quad C_r=2rC_0\rho^{-(r-1)},\quad
\bar C=C_0\left(1+\frac{D}{e\alpha\rho\log(1/\rho)}\right).
\]
If $\ell\ge\max(1,r-1)$ and $C_r(1+D\ell/\rho)\rho^{\alpha\ell}\le1$, then
\begin{equation}
\norm{\delta y_t}\le\bar C\bar\rho^t\norm{\delta z_0}_{\mathcal Z}
+\bar CC_r\sum_{i=0}^{t-1}\bar\rho^{t-1-i}\norm{\delta v_i}.
\label{eq:chunkbound}
\end{equation}
\emph{Proof.} Put $Z_j=\norm{\delta z_{j\ell}}_{\mathcal Z}$ and $V_i=\norm{\delta v_i}$. The policy gives $\norm{\delta u_{j\ell+k}}\le D\rho^kZ_j+V_{j\ell+k}$. Substituting in \eqref{eq:ios} yields, for $0\le k\le\ell$,
\[
\norm{\delta y_{j\ell+k}}\le C_0(1+Dk/\rho)\rho^kZ_j
+C_0\sum_{i=0}^{k-1}\rho^{k-1-i}V_{j\ell+i}.
\]
Since $\ell\ge r-1$, all blocks of the next history lie within the chunk, including its initial observation when equality holds. Summing their bounds gives
\[
Z_{j+1}\le C_r(1+D\ell/\rho)\rho^\ell Z_j
+C_r\sum_{i=0}^{\ell-1}\rho^{\ell-1-i}V_{j\ell+i}.
\]
The first coefficient is at most $\bar\rho^\ell$. Iterate, use $\rho\le\bar\rho$, and apply $k\rho^{\alpha k}\le1/[e\alpha\log(1/\rho)]$ to the within-chunk bound to obtain \eqref{eq:chunkbound}. $\square$

Sufficiently long chunks satisfy the condition when $H$ permits them. However, decreasing chunk sensitivity is an assumption, and the bound becomes more conservative as $r$ grows. It does not by itself prove that context permits shorter execution.

\paragraph{Exact linear test and example.}
Augment the state with observation/action buffers: $s_{t+1}=\mathsf As_t+\mathsf Bu_t+\mathsf Jw_t$, $z_t=\mathsf Ps_t$. A fitted linear chunk policy has boundary map
\begin{equation}
M_{r,\ell}=\mathsf A^\ell+\sum_{k=0}^{\ell-1}\mathsf A^{\ell-1-k}\mathsf B\widehat W_k\mathsf P.
\label{eq:lift}
\end{equation}
Unrolling gives the map. The augmented system is exponentially stable exactly when $\spr(M_{r,\ell})<1$. For $x_+=x/2+u$ and expert $K=0$, the imperfect predictor $\pi_k(x)=3\cdot2^{-k}x/4$ gives $M_\ell=2^{-\ell}(1+3\ell/2)$. At $\ell=1,2,3$, the multipliers are $1.25,1,0.6875$. Three-step commitment stabilizes this predictor family; the exact expert still needs no chunking.

\clearpage
\section{Supporting evidence and remaining objective}
The main simulations use stable $A=0.995P_n$, $B=b$, $K=-0.05b^\top P_n$, where $P_n$ is a cyclic permutation and $b$ a unit vector. Ridge fits use 32 noiseless demonstrations, 24 matched windows each, penalty $10^{-6}$, and $H=16$. Three seeds each supply eight paired episodes with shared initial states and process noise of coordinate standard deviation $0.02$. After a common $n$-step expert warm-up, tail error averages $\norm{x_t-x_t^*}^2/n$ over the last 200 of 600 steps. Output-separable regression makes the first-action fit identical across prediction horizons, isolating the effect of execution.

\begin{center}\small
\begin{tabular}{ccrr}\toprule
Context $r$ & Execution $\ell$ & Mean tail MSE & Max. per-step rate\\\midrule
3&1&$\infty$&1.3215\\
3&2&$2.80\times10^{-4}$&0.9923\\
4&1&$7.64\times10^{-6}$&0.9923\\
4&16&$5.55\times10^{-5}$&0.9923\\\bottomrule
\end{tabular}
\end{center}
For $n=10,h=4$, the table reports rates $\spr(M_{r,\ell})^{1/\ell}$. At $r=3$, validation MSE is $4.62\times10^{-14}$ despite failure in all 24 single-step rollouts. Two-step commitment stabilizes the fit; context four enables single-step execution with $37\times$ lower tracking error. Across 20 dimension/sensor configurations, single-step execution is best in mean tail error in 19 cases at $r=h$. These observations illustrate the mechanisms, not Proposition 3's policy-class assumptions.

We also study the challenging pair of \citet{simchowitz2025}. No narrow-data setting meets the criterion of stability and tracking RMS at most $0.05$ for all five seeds, whereas persistent excitation with full observations permits one-step execution at every tested context. History alone does not resolve identification. The remaining theoretical task is to bound deployment error or feedback gain as context grows, then characterize the shortest reliable execution. Observer-based baselines, noise-conditioning analysis, and joint learner--expert covariance calculations will connect that theory to further simulations. The current finite-horizon results concern fixed plants and few seeds; they do not establish a universal context--execution tradeoff. Stability alone is also insufficient: zero control already stabilizes a Schur-stable plant, so imitation must be assessed through expert tracking.

\begin{thebibliography}{2}
\bibitem[Simchowitz et al.(2025)]{simchowitz2025}
M. Simchowitz, D. Pfrommer, and A. Jadbabaie.
The pitfalls of imitation learning when actions are continuous.
\emph{arXiv:2503.09722}, 2025.
\bibitem[Zhang et al.(2025)]{zhang2025}
T. T. Zhang, D. Pfrommer, C. Pan, N. Matni, and M. Simchowitz.
Action chunking and exploratory data collection yield exponential improvements in behavior cloning for continuous control.
\emph{arXiv:2507.09061}, 2025.
\end{thebibliography}
\end{document}
